\documentclass[10pt,a4paper]{article}
\usepackage[top=18mm,bottom=18mm,left=17mm,right=17mm]{geometry}
\usepackage[T1]{fontenc}
\usepackage{lmodern,graphicx,tabularx,array,booktabs,xcolor,hyperref,fancyhdr}
\hypersetup{colorlinks=true,urlcolor=blue,linkcolor=blue,pdftitle={DES-001 RD53i module proposal: reviewer brief},pdfauthor={nODD project}}
\pdfinfoomitdate=1
\pdftrailerid{}
\setlength{\parindent}{0pt}
\setlength{\parskip}{5pt}
\setlength{\tabcolsep}{4pt}
\renewcommand{\arraystretch}{1.18}
\pagestyle{fancy}
\fancyhf{}
\fancyhead[L]{\small nODD \quad DES-001 reviewer brief}
\fancyhead[R]{\small DRAFT \quad 2026-09-17}
\fancyfoot[L]{\small Module-only proposal; human review pending}
\fancyfoot[R]{\small \thepage/3}
\setlength{\headheight}{14pt}
\newcommand{\repo}[2]{\href{https://github.com/asalzburger/nodd/blob/design/rd53-pixel-modules/#1}{#2}}
\newcommand{\heading}[1]{\vspace{3pt}{\large\bfseries #1}\par}
\newcolumntype{Y}{>{\raggedright\arraybackslash}X}
\begin{document}
{\LARGE\bfseries RD53i pixel modules}\par
{\large First-review proposal summary}\par
\textbf{Recommendation.} Qualify the single-chip A1 assembly first, then compare
the two-chip A2 and four-chip B4 assemblies before choosing a coverage module.
Use the fewest module types justified by evidence; A1 may remain a study unit.
Neither a complete material budget nor a coverage winner is established.

\heading{What is fixed, and what this review covers}
RD53i is the selected chip, with its family and revision already resolved;
terminology is defined once in the
\repo{docs/design/DES-001-rd53-pixel-modules.md\#glossary}{DES-001 glossary}.
The design principles are reusable structures, few module types, low accounted
material and credible assembly. Independent Astra agents rewrote both
alternatives; a separate Astra Coordinator assessed their tradeoffs.

The assembly includes the planar sensor, chip substrates, bump interconnects,
flex and adhesives, wire bonds, sensor-bias contact, local passives and
module electrical termination. Any integral flex extension is part of the BOM.
The boundary ends at that termination and bare chip backs.
\textbf{Mounting structures, mounting adhesive, cooling infrastructure and
external power/data distribution are excluded.} Pixel volume and placement
remain deliberately open. No detector geometry is implemented.

\heading{The alternatives at a glance}
\begin{tabularx}{\linewidth}{@{}lYYY@{}}
\toprule
 & \textbf{A1 single chip} & \textbf{A2 two chips} & \textbf{B4 quad}\\
\midrule
Assembly & One sensor, chip and local flex & Two side-by-side chips under one sensor/flex & Four chips in a $2\times2$ sensor/flex assembly\\
Nominal channels & 153,600 & 307,200 & 614,400\\
Readout footprint & 384 mm$^2$ & 768 mm$^2$ & 1536 mm$^2$\\
Approx. die-only span & $20\times21$ mm & $40\times21$ mm & $40\times42$ mm\\
Main opportunity & Small qualification and rejected-area unit & Shares local components without a quad-sized assembly & Amortizes repeated module edges and termination\\
Main challenge & More repeated edges/terminations per area & Separate sensor/flex and seam qualification & Larger shared sensor, more routing and correlated assembly losses\\
\bottomrule
\end{tabularx}
\textbf{Read these numbers as bookkeeping.} They derive from the chip overview's
$400\times384$ cells at 50 $\mu$m pitch and approximate die size (PM3-F01,
PM3-I01). Die-only spans exclude gaps, guard regions, bond access, flex and
termination. Readout footprint is not measured efficient area or a finished
sensor/module outline. A shared sensor does not establish seamless response.

\heading{What the reviewer is being asked to assess}
Are the proposed hybrid topology and assembly sequence credible? Are the
nominal thicknesses and component inventory a sound starting point? Does the
A1 qualification-first programme collect the right evidence to compare A2 and
B4 while avoiding an unjustified three-type family?
\textbf{This is first technical review, not full design sign-off.}
DES-001 and ADR-007 remain DRAFT; reviewers and numeric acceptance criteria
are not yet assigned.

\newpage
{\LARGE\bfseries Material study and physical stack}\par
\textbf{Proposed thicknesses; sourced radiation lengths; derived contributions.}
The same nominal layer stack controls the comparison. These thicknesses are
nODD study choices (PM3-C02--C07), not verified delivered parts. Material
$X_0$ constants are public-source facts (PM3-F02--F05). Percentages are
normal-crossing calculations (PM3-I02).

{\small
\begin{tabularx}{\linewidth}{@{}p{35mm}p{32mm}Yrr@{}}
\toprule
\textbf{Component} & \textbf{Thickness / geometry} & \textbf{Material} & \textbf{$X_0$ [mm]} & \textbf{\% $X_0$}\\
\midrule
Sensor bulk & 150 $\mu$m & Silicon & 93.70 & 0.1601\\
Chip substrate & 150 $\mu$m per die & Silicon & 93.70 & 0.1601\\
Flex core + coverlay & $25+2\times12.5=50$ $\mu$m & Specified polyimide film & 285.7 & 0.0175\\
Copper routing & $2\times18=36$ $\mu$m & Copper, both layers crossed & 14.36 & 0.2507\\
Flex-to-sensor bond & 25 $\mu$m nominal & Epoxy grade/filler TBD & TBD & TBD\\
Bump interconnect & 20 $\mu$m standoff & Alloy/under-bump stack TBD & TBD & TBD\\
Wire bonds & 25 $\mu$m diameter; loops TBD & Aluminum candidate & 88.97 & TBD\\
Laminate adhesives & Thickness/coverage TBD & Supplier grades TBD & TBD & TBD\\
Device films/finishes & Thickness/coverage TBD & Metals/passivation/plating & TBD & TBD\\
Local parts/termination & Part volumes/masses TBD & Mixed spatial BOM & TBD & TBD\\
\bottomrule
\end{tabularx}}

\textbf{Partial local reference subtotal: 0.5884\% $X_0$.}
For the four calculable rows, $100\sum t_i/X_{0,i}$ uses matched units and
unrounded inputs. It crosses the sensor, \emph{one} chip, all polyimide films and
both copper layers. The four quad chips are lateral, not four stacked slabs.

This is \textbf{not a complete module budget, an area average or a bound on the
completed module}. Actual routed copper coverage changes the spatial result.
Glue, bumps, laminate adhesives, films, wire loops and local parts are additional
unresolved contributions, not zero. Bump standoff and wire diameter are not
area-equivalent metal thicknesses. Integral flex extensions and termination
must be counted, even outside the sensor projection.

\heading{Cross-section at an accessible chip edge}
\par\noindent\includegraphics[width=\linewidth]{cross-section.pdf}\par
{\small Crop of the original compact schematic. Layer order is shared by both
alternatives: sensor-backside flex, adhesive, sensor, discrete bumps and chip.
Bond loops go around the sensor edge to accessible flex lands; exact drawings
must establish pad access. Display heights are schematic; thickness labels are
proposed study inputs. Unknown adhesive/film layers remain additional.}

\textbf{Why equal subtotals do not rank the modules.} Material per useful area
depends on sensor margins/seams, routed flex, local parts, bonds, termination and
assembly yield. Complete each spatial BOM before claiming a saving.

\newpage
{\LARGE\bfseries Module layouts and review priorities}\par
\textbf{A1 / optional A2:} common chip/interface principles, but A2 requires a
qualified shared sensor, seam mapping and its own routed flex. Keep A1 as the
smallest learning unit; justify any retained production variant.
\par\noindent\includegraphics[width=\linewidth]{compact-plan.pdf}\par

\textbf{B4:} four chips share a larger sensor and flex. Outward chip peripheries
preserve proposed bond access; both seams and their intersection need response
and assembly evidence. The A--A line identifies the opposing-chip section in the
full quad drawing.
\par\noindent\includegraphics[width=\linewidth]{quad-plan.pdf}\par
{\small Original plan-view crops; schematic, not to scale. Colors identify
components within each drawing. Exact sensor edges, gaps, loops and routing
remain unresolved; nominal readout rectangles do not certify acceptance.}

\heading{Priorities before a family recommendation}
\begin{enumerate}
\setlength{\itemsep}{1pt}
\item Obtain authoritative die/pad/bump and electrical interfaces, delivered
thinning and known-good-die evidence; confirm bond access and sensor compatibility.
\item Define sensor process/bias and operating requirements; dimension seams,
guard regions, flex, bonds and termination without hiding inactive regions.
\item Complete component grades, thicknesses, volumes and spatial material
maps; qualify module connectivity, data integrity, metrology and seam response.
\item Compare accepted-area yield, rework and partial-operation behavior.
Choose the minimum justified family; defer deployment suitability to pixel-volume work.
\end{enumerate}

{\small\textbf{Source record and detailed proposals.}
\repo{docs/design/DES-001-rd53-pixel-modules.md}{DES-001 (PM3 claims and glossary)};
\repo{docs/design/inputs/DES-001-compact-modules.md}{compact A};
\repo{docs/design/inputs/DES-001-quad-modules.md}{quad B};
\repo{docs/design/inputs/DES-001-coordinator-review.md}{Coordinator assessment};
\repo{docs/design/DES-001-module-stack.json}{common inputs/unrounded arithmetic};
\repo{reference/manifest.yaml}{source catalogue}.
Radiation-length rows: PDG Si/Cu/Al 2025 and polyimide film 2020.
This brief summarizes existing claims; it adds no approved parameters or hardware validation.}
\end{document}
